
\documentclass[aspectratio=169]{beamer}
\usepackage[spanish]{babel}
\usepackage[utf8]{inputenc}
\usepackage[T1]{fontenc}
\usepackage{siunitx}
\usepackage{amsmath,amssymb}
\usetheme{Madrid}

\title[Criterios de modelo ideal]{Criterios para la elección del modelo ideal\\en flujo compresible en conductos}
\author{}
\date{}

\begin{document}

%------------------------------------------------
\begin{frame}
  \titlepage
\end{frame}
%------------------------------------------------

\begin{frame}{Punto de partida: expresión general para $w/A$}
Del balance de energía mecánica en una conducción horizontal, sección constante y flujo estacionario compresible:

\begin{equation*}
\left(\frac{w}{A}\right)^2 \ln\!\frac{v_2}{v_1}
\;+\;\int_1^2 \frac{\mathrm dP}{v}
\;+\;2 f \frac{L}{D}\left(\frac{w}{A}\right)^2 = 0
\end{equation*}

\pause

\begin{equation*}
\frac{w}{A}
=
\left[
-\frac{\displaystyle \int_1^2 \frac{\mathrm dP}{v}}
{\displaystyle \ln\!\frac{v_2}{v_1} + 2 f\frac{L}{D}}
\right]^{1/2}.
\end{equation*}

\pause

\begin{itemize}
  \item Para comparar modelos (isotérmico, adiabático, con calor, etc.) mantenemos:
  \begin{itemize}
    \item misma caída de presión $\Delta P$,
    \item mismo término de fricción $2fL/D$,
  \end{itemize}
  \item y analizamos cómo cambian $-\int dP/v$ y $\ln(v_2/v_1)$.
\end{itemize}
\end{frame}

%------------------------------------------------

\begin{frame}{Isotérmico vs.\ adiabático (mismos $P_1,P_2$)}
Gas ideal: $Pv/T = \text{cte}$.

\begin{columns}[T]
\column{0.48\textwidth}
\textbf{Isotérmico}
\begin{itemize}
  \item $T_2 = T_1$.
  \item $v_{2,i}$ mayor.
\end{itemize}

\column{0.48\textwidth}
\textbf{Adiabático}
\begin{itemize}
  \item $T_2 < T_1$.
  \item $v_{2,a}$ menor.
\end{itemize}
\end{columns}

\medskip
Para el mismo $(P_1,P_2)$:

\begin{align*}
-\int_1^2 \frac{\mathrm dP}{v}\bigg|_i
&<
-\int_1^2 \frac{\mathrm dP}{v}\bigg|_a,\\[4pt]
\ln\frac{v_{2,i}}{v_1}
&>
\ln\frac{v_{2,a}}{v_1}.
\end{align*}

\pause

\begin{block}{Conclusión}
\[
\boxed{
\left(\frac{w}{A}\right)_{\text{iso}}
<
\left(\frac{w}{A}\right)_{\text{adiab}}
}
\]

El modelo adiabático predice \emph{mayor} flujo que el isotérmico para la misma $\Delta P$.
\end{block}
\end{frame}

%------------------------------------------------

\begin{frame}{Efecto del calor respecto al adiabático}
Comparamos con el caso adiabático ($Q=0$), mismo par $(P_1,P_2)$:

\begin{itemize}
  \item \textbf{Calentamiento} ($Q>0$):
  \begin{itemize}
    \item $T$ y $v$ mayores que en el adiabático.
    \item $-\displaystyle\int dP/v\big|_{Q>0}
           < -\displaystyle\int dP/v\big|_{a}$.
  \end{itemize}
  \item \textbf{Enfriamiento} ($Q<0$):
  \begin{itemize}
    \item $v$ menor que en el adiabático.
    \item $-\displaystyle\int dP/v\big|_{Q<0}
           > -\displaystyle\int dP/v\big|_{a}$.
  \end{itemize}
\end{itemize}

\pause

\begin{block}{Orden de magnitud de $w/A$}
\[
\boxed{
\left(\frac{w}{A}\right)_{Q>0}
<
\left(\frac{w}{A}\right)_{\text{adiab}}
<
\left(\frac{w}{A}\right)_{Q<0}
}
\]

Calentar el gas reduce el flujo respecto al adiabático; enfriarlo lo aumenta.
\end{block}
\end{frame}

%------------------------------------------------

\begin{frame}{Perfil de temperatura vs.\ isotérmico a $T_1$}
Tomamos como referencia el modelo isotérmico a $T_1$.

\begin{itemize}
  \item Si en toda la tubería $T > T_1$:
  \[
    -\int \frac{\mathrm dP}{v}\bigg|_{T>T_1}
    <
    -\int \frac{\mathrm dP}{v}\bigg|_{T=T_1}.
  \]
  \item Si en toda la tubería $T < T_1$:
  \[
    -\int \frac{\mathrm dP}{v}\bigg|_{T<T_1}
    >
    -\int \frac{\mathrm dP}{v}\bigg|_{T=T_1}.
  \]
\end{itemize}

\pause

\begin{block}{Orden de $w/A$ respecto al isotérmico a $T_1$}
\[
\boxed{
\left(\frac{w}{A}\right)_{T>T_1}
<
\left(\frac{w}{A}\right)_{\text{iso a }T_1}
<
\left(\frac{w}{A}\right)_{T<T_1}
}
\]

Si el gas está más caliente que $T_1$ a lo largo de la línea, el flujo real
es menor que el del modelo isotérmico a $T_1$; si está más frío, es mayor.
\end{block}
\end{frame}

%------------------------------------------------

\begin{frame}{Mapa cualitativo: casos típicos (I)}
Ubicamos $w_{\text{real}}$ entre $w_i$ (isotérmico) y $w_a$ (adiabático) usando $T_1$, $T_{\text{amb}}$ y el signo de $Q$.

\medskip

\textbf{Caso 1: gas muy frío, ambiente caliente}
\begin{itemize}
  \item $T_1 \ll T_{\text{amb}}$, $Q>0$, con $T(x) > T_1$.
  \item El gas se calienta fuertemente.
\end{itemize}

\[
w_{\text{real}} < w_i < w_a.
\]

\pause

\textbf{Caso 2: gas a temperatura similar a la ambiente}
\begin{itemize}
  \item $T_1 \approx T_{\text{amb}}$, $Q>0$ moderado,
        $T(x)$ apenas por debajo de $T_1$.
\end{itemize}

\[
w_i < w_{\text{real}} < w_a.
\]

El flujo real queda entre ambos modelos; el isotérmico es cota inferior, el adiabático cota superior.
\end{frame}

%------------------------------------------------

\begin{frame}{Mapa cualitativo: casos típicos (II)}
\textbf{Caso 3: gas muy caliente, ambiente frío (enfriamiento intenso)}
\begin{itemize}
  \item $T_1 \gg T_{\text{amb}}$, $Q<0$, con $T(x)<T_1$.
  \item El gas se enfría fuertemente a lo largo de la conducción.
\end{itemize}

\[
w_i < w_a < w_{\text{real}}.
\]

\pause

\textbf{Casos intermedios}
\begin{itemize}
  \item Si $T_1 < T_{\text{amb}}$ y hay calentamiento pero no muy fuerte:
  \[
  w_{\text{real}} < w_a.
  \]
  \item Si $T_1 > T_{\text{amb}}$ y el gas se mantiene más frío que $T_1$:
  \[
  w_i < w_{\text{real}}.
  \]
\end{itemize}
\end{frame}

%------------------------------------------------

\begin{frame}{Resumen de desigualdades clave}
\begin{block}{Isotérmico vs.\ adiabático}
\[
\left(\frac{w}{A}\right)_{\text{iso}}
<
\left(\frac{w}{A}\right)_{\text{adiab}}.
\]
\end{block}

\pause

\begin{block}{Con calor vs.\ adiabático}
\[
\left(\frac{w}{A}\right)_{Q>0}
<
\left(\frac{w}{A}\right)_{\text{adiab}}
<
\left(\frac{w}{A}\right)_{Q<0}.
\]
\end{block}

\pause

\begin{block}{Perfil térmico vs.\ isotérmico a $T_1$}
\[
\left(\frac{w}{A}\right)_{T>T_1}
<
\left(\frac{w}{A}\right)_{\text{iso a }T_1}
<
\left(\frac{w}{A}\right)_{T<T_1}.
\]
\end{block}

Estas relaciones sirven como ``faro'' para ubicar $w_{\text{real}}$ entre los modelos ideales.
\end{frame}

%------------------------------------------------

\begin{frame}{Criterios prácticos de elección de modelo}
\textbf{1.\ Modelo adiabático}
\begin{itemize}
  \item Conducción bien aislada o tiempo de residencia corto.
  \item En general, lo uso como \textbf{cota superior} de flujo para una misma $\Delta P$.
  \item Excepción: enfriamientos intensos ($Q<0$ grande) que pueden llevar a $w_{\text{real}} > w_a$.
\end{itemize}

\pause

\textbf{2.\ Modelo isotérmico (a $T_1$ o $T_{\text{amb}}$)}
\begin{itemize}
  \item Tubería bien acoplada térmicamente al ambiente,
        $T_1 \simeq T_{\text{amb}}$.
  \item En general, \textbf{cota inferior} de flujo para la misma $\Delta P$.
\end{itemize}
\end{frame}

%------------------------------------------------

\begin{frame}{Ubicación cualitativa del caso real}
\textbf{Reglas de pulgar}:

\begin{itemize}
  \item Si $Q>0$ (calentamiento global):
  \[
  w_{\text{real}} \le w_{\text{adiab}},
  \quad
  \text{frecuentemente } w_{\text{real}} \lesssim w_{\text{iso}}.
  \]
  \item Si $Q<0$ (enfriamiento global):
  \[
  w_{\text{real}} \ge w_{\text{adiab}},
  \quad
  \text{típicamente } w_{\text{real}} \gtrsim w_{\text{iso}}.
  \]
  \item La relación $T_1$--$T_{\text{amb}}$ indica si $w_{\text{real}}$ queda:
  \begin{itemize}
    \item por debajo de ambos modelos,
    \item entre ambos,
    \item o por encima de ambos,
  \end{itemize}
  siguiendo los casos 1--3 del mapa cualitativo.
\end{itemize}

\bigskip
En resumen, uso isotérmico y adiabático como modelos ideales ``faro'' y ubico
el caso real entre ambos para justificar la elección del modelo en cada problema.
\end{frame}

%------------------------------------------------

\end{document}
